\hypertarget{pypy-the-jit-and-meta-tracing}{%
\chapter{PyPy: The JIT and
Meta-Tracing}\label{pypy-the-jit-and-meta-tracing}}

PyPy is an alternative implementation of Python written in
\textbf{RPython} (a restricted subset of Python).

\hypertarget{the-meta-tracing-jit}{%
\subsection{74.1 The Meta-Tracing JIT}\label{the-meta-tracing-jit}}

Unlike the CPython JIT (Chapter 21), which compiles bytecode to machine
code, PyPy's JIT ``traces'' the execution of the interpreter itself. *
\textbf{Warming Up}: The first time a loop runs, it is interpreted. *
\textbf{Recording}: If the loop is ``hot,'' the JIT records every
instruction. * \textbf{Optimizing}: It removes redundant operations
(like repeated type checks for the same object) and generates highly
optimized machine code.

\hypertarget{memory-management-the-incrgc}{%
\subsection{74.2 Memory Management: The
IncrGC}\label{memory-management-the-incrgc}}

PyPy uses a moving, generational garbage collector. This is more
efficient than reference counting for long-running processes but can
lead to higher ``pause times'' during major collections.

\begin{center}\rule{0.5\linewidth}{0.5pt}\end{center}

\hypertarget{appendices}{%
\chapter{Appendices}\label{appendices}}
